\section{Introduction}

Zero-shot variant effect predictors (VEPs), from protein language models such as ESM-1b~\cite{rives2021biological,brandes2023genome} to alignment-based models such as EVE and GEMME~\cite{frazer2021disease,laine2019gemme} and hybrids such as TranceptEVE and PoET~\cite{notin2022trancepteve,truong2023poet}, now rank missense variants with high accuracy on clinical benchmarks~\cite{notin2023proteingym}. Clinical use, however, needs more than a ranking: variant interpretation under the guidelines of the American College of Medical Genetics and Genomics and the Association for Molecular Pathology (ACMG/AMP)~\cite{richards2015standards} needs evidence whose error rate is known, and the ClinGen calibration of computational evidence~\cite{pejaver2022calibration} has shown that a single genome-wide calibration can hide large differences between genes~\cite{tejura2024calibration}.

Conformal prediction~\cite{vovk2005algorithmic,angelopoulos2023gentle} offers distribution-free error control: a prediction set over $\{\text{benign},\text{pathogenic}\}$ that contains the true label with probability at least $1-\alpha$. Its guarantee rests on exchangeability between calibration and test data. For VEPs this assumption fails in two ways. First, variants are grouped into proteins and homologous protein families: a predictor's accuracy differs from family to family, and families contribute very different numbers of labeled variants (from \minVar{} to \maxVar{} per protein in our data). Pooling variants lets the largest families set the error rate, so the ``90\%'' is a statement about variants from well-studied genes, not about a new protein. Second, the proteins a predictor is deployed on are often evolutionarily distant from, or have shallower alignments than, the proteins in the calibration set: a family-level distribution shift that is visible without labels.

We make three contributions.
\begin{itemize}
\item We formalize \emph{family-weighted} coverage, the probability that a random variant of a new protein family is covered, and its class-conditional version for pathogenic variants. On ProteinGym clinical substitutions, class-conditional (Mondrian) split conformal, the standard remedy for class imbalance, meets its own variant-level target yet covers only \sOneBtwoPath{} of pathogenic variants per family for four of five predictors at a 0.90 target, with no distribution shift at all.
\item We present ProtCOVER, which (i) calibrates at the family level, (ii) reweights calibration families by a cross-fitted density ratio of label-free family descriptors, and (iii) gates queries with insufficient calibration support. We prove finite-sample family-weighted coverage under family exchangeability (Proposition~\ref{prop:hcp}, a restatement of hierarchical conformal prediction~\cite{lee2026hierarchical}) and under family-level covariate shift (Theorem~\ref{thm:weighted}), bound the loss from estimated weights (Corollary~\ref{cor:estimated}), and show that the gate preserves all of these (Lemma~\ref{lem:enlarge}).
\item We evaluate five predictors under exchangeable, synthetic-shift, and orphan-family splits with family-level bootstrap intervals. Under shift, ProtCOVER keeps both class-conditional coverages at or above \worstPw{} for every predictor and setting, against \worstHcp{} for family-level calibration alone and \worstGwcp{} for group-weighted conformal prediction; the gain is largest when the shift acts on the descriptors it models, it matches group-weighted conformal prediction on orphan families, and it costs larger sets. We also report a negative result: abstaining on poorly supported families did not reduce the error of the remaining predictions relative to random abstention.
\end{itemize}
